\documentclass[11pt,a4paper]{article}
\usepackage[a4paper,margin=24mm]{geometry}
\usepackage[utf8]{inputenc}
\usepackage[T1]{fontenc}
\usepackage{lmodern}
\usepackage{amsmath,amssymb,amsthm}
\usepackage{booktabs,array,tabularx}
\usepackage{enumitem}
\usepackage{float}
\usepackage{longtable}
\usepackage{graphicx}
\usepackage{microtype}
\usepackage{xcolor}
\definecolor{linkblue}{rgb}{0,0.15,0.5}
\usepackage[colorlinks=true,linkcolor=linkblue,citecolor=linkblue,urlcolor=linkblue]{hyperref}
\usepackage{url}
\allowdisplaybreaks
\emergencystretch=2em

\renewcommand{\thesection}{S\arabic{section}}
\renewcommand{\theequation}{S.\arabic{equation}}
\renewcommand{\thetable}{S\arabic{section}.\arabic{table}}
\renewcommand{\thefigure}{S\arabic{section}.\arabic{figure}}
\theoremstyle{plain}
\newtheorem{theorem}{Theorem}[section]
\newtheorem{proposition}[theorem]{Proposition}
\newtheorem{lemma}[theorem]{Lemma}
\newtheorem{corollary}[theorem]{Corollary}
\theoremstyle{definition}
\newtheorem{definition}[theorem]{Definition}
\theoremstyle{remark}
\newtheorem{remark}[theorem]{Remark}

\newcommand{\R}{\mathbb{R}}
\newcommand{\Qr}{\mathbb{Q}}
\newcommand{\N}{\mathbb{N}}
\newcommand{\calA}{\mathcal{A}}
\newcommand{\calF}{\mathcal{F}}
\newcommand{\calK}{\mathcal{K}}
\newcommand{\calN}{\mathcal{N}}
\newcommand{\calP}{\mathcal{P}}
\newcommand{\calS}{\mathcal{S}}
\newcommand{\calW}{\mathcal{W}}
\newcommand{\Vstar}{V^{\ast}}
\newcommand{\Rstar}{R^{\ast}}
\newcommand{\Gstar}{G^{\ast}}
\newcommand{\tfr}[2]{\tfrac{#1}{#2}}
\newcommand{\st}[1]{\mathsf{#1}}
\newcommand{\led}{\ell}
\newcommand{\priv}{\pi}
\newcommand{\opt}{\ell^{+}}
\newcommand{\free}{\mathrm{f}}
\newcommand{\nr}{\mathrm{nr}}
\newcommand{\ofull}{o_{\mathrm{id}}}
\newcommand{\oagg}{o_{\Sigma}}
\newcommand{\oblind}{o_{\mathrm{b}}}
\newcommand{\Kinst}{K_{\mathrm{inst}}}
\newcommand{\Kagg}{K_{\mathrm{agg}}}
\newcommand{\Kcod}{K_{\mathrm{cod}}}
\newcommand{\Kfull}{K_{\mathrm{full}}}
\newcommand{\Kdec}{K_{\mathrm{dec}}}
\newcommand{\NCHECKS}{64}
\newcommand{\mref}[1]{\csname mainlab@#1\endcsname}
\newcommand{\meqref}[1]{(\mref{#1})}
\setlength{\parskip}{2pt}
\renewcommand{\arraystretch}{1.2}

\begin{document}

\title{\textbf{Supplementary material for}\\ \emph{Worked systems for viability under incomplete observation: exact belief kernels across observation structures, protocols and decentralisation, a monitoring-design census, and certified margins}}
\author{Amin Abaee\\[0.3em]{\small Independent Researcher}\\[0.2em]{\small ORCID: \href{https://orcid.org/0000-0002-0019-1842}{0000-0002-0019-1842}\quad \href{mailto:amin_abaee@ut.ac.ir}{amin\_abaee@ut.ac.ir}}}
\date{October 15, 2026}
\maketitle

\noindent This supplement collects the tables behind the counts of the main text and the material removed from it: the master table of the $36$ pairs (S1), the whole-universe census with the register conventions that are not used in the main text (S2), the stationary-policy enumeration (S3), the decentralised law tables (S4), the monitoring-design census (S5), the four-state target census (S6), the timing grids (S7), the scaling of the census to larger caps (S8), the continuous-fibre tables (S9), the verifier output and the Lean declarations (S10), and a table of symbols (S11). Theorem, section and equation numbers of the form ``Theorem~\mref{thm:mono}'' refer to the main text. Every table is generated from \texttt{results.json}.

% generated from the main-text .aux file: main-text numbers used by \mref / \meqref
\expandafter\def\csname mainlab@sec:intro\endcsname{1}
\expandafter\def\csname mainlab@sec:frame\endcsname{2}
\expandafter\def\csname mainlab@sec:system\endcsname{2.1}
\expandafter\def\csname mainlab@def:system\endcsname{2.1}
\expandafter\def\csname mainlab@eq:R\endcsname{1}
\expandafter\def\csname mainlab@fig:system\endcsname{1}
\expandafter\def\csname mainlab@def:obs\endcsname{2.2}
\expandafter\def\csname mainlab@def:protocol\endcsname{2.3}
\expandafter\def\csname mainlab@sec:game\endcsname{2.2}
\expandafter\def\csname mainlab@def:game\endcsname{2.4}
\expandafter\def\csname mainlab@eq:Phi\endcsname{2}
\expandafter\def\csname mainlab@lem:basic\endcsname{2.5}
\expandafter\def\csname mainlab@def:classes\endcsname{2.6}
\expandafter\def\csname mainlab@sec:mono\endcsname{3}
\expandafter\def\csname mainlab@thm:mono\endcsname{3.1}
\expandafter\def\csname mainlab@cor:order\endcsname{3.2}
\expandafter\def\csname mainlab@prop:strict\endcsname{3.3}
\expandafter\def\csname mainlab@rem:lean\endcsname{3.4}
\expandafter\def\csname mainlab@sec:census\endcsname{4}
\expandafter\def\csname mainlab@sec:census_numbers\endcsname{4.1}
\expandafter\def\csname mainlab@prop:full\endcsname{4.1}
\expandafter\def\csname mainlab@tab:kernels\endcsname{1}
\expandafter\def\csname mainlab@fig:lattice\endcsname{2}
\expandafter\def\csname mainlab@sec:alt\endcsname{4.2}
\expandafter\def\csname mainlab@lem:alt\endcsname{4.2}
\expandafter\def\csname mainlab@thm:closed\endcsname{4.3}
\expandafter\def\csname mainlab@cor:conv\endcsname{4.4}
\expandafter\def\csname mainlab@sec:stationary\endcsname{4.3}
\expandafter\def\csname mainlab@prop:stationary\endcsname{4.5}
\expandafter\def\csname mainlab@sec:dec\endcsname{5}
\expandafter\def\csname mainlab@def:laws\endcsname{5.1}
\expandafter\def\csname mainlab@thm:factor\endcsname{5.2}
\expandafter\def\csname mainlab@prop:deccensus\endcsname{5.3}
\expandafter\def\csname mainlab@sec:monitor\endcsname{6}
\expandafter\def\csname mainlab@def:adequacy\endcsname{6.1}
\expandafter\def\csname mainlab@thm:adequacy\endcsname{6.2}
\expandafter\def\csname mainlab@eq:count\endcsname{3}
\expandafter\def\csname mainlab@thm:dynamic\endcsname{6.3}
\expandafter\def\csname mainlab@thm:censusP\endcsname{6.4}
\expandafter\def\csname mainlab@fig:census\endcsname{3}
\expandafter\def\csname mainlab@tab:design\endcsname{2}
\expandafter\def\csname mainlab@sec:timing\endcsname{7}
\expandafter\def\csname mainlab@def:regime\endcsname{7.1}
\expandafter\def\csname mainlab@prop:regime\endcsname{7.2}
\expandafter\def\csname mainlab@fig:timing\endcsname{4}
\expandafter\def\csname mainlab@prop:maxplus\endcsname{7.3}
\expandafter\def\csname mainlab@cor:grid\endcsname{7.4}
\expandafter\def\csname mainlab@sec:continuous\endcsname{8}
\expandafter\def\csname mainlab@prop:bench\endcsname{8.1}
\expandafter\def\csname mainlab@prop:duality\endcsname{8.2}
\expandafter\def\csname mainlab@prop:helly\endcsname{8.3}
\expandafter\def\csname mainlab@fig:continuous\endcsname{5}
\expandafter\def\csname mainlab@prop:drift\endcsname{8.4}
\expandafter\def\csname mainlab@sec:rules\endcsname{9}
\expandafter\def\csname mainlab@sec:verification\endcsname{10}
\expandafter\def\csname mainlab@sec:related\endcsname{11}
\expandafter\def\csname mainlab@sec:conclusion\endcsname{12}
\section{The master table of the 36 pairs}\label{si:master}

The earlier editions of the audit worked on the $36$ unordered pairs of safe states; Table~\ref{tab:master} retains this population and extends it with the private-register and optimistic-ledger conventions and the decentralised kernel. The eight pairs through the corner $\st{11}$ are dead in every class because the corner has no safe instrument (Lemma~\mref{lem:basic}(c)). Among the $28$ corner-free pairs the dead ones are: institutional ($\led$) $\st{12|21},\st{12|31},\st{13|21},\st{13|31}$; aggregate $\st{12|21},\st{13|31}$; codex ($\led$) $\st{12|21},\st{12|31},\st{13|21}$; full none. The mechanisms are as follows.
\begin{itemize}[leftmargin=2em,itemsep=1pt]
\item $\st{12|21}$ under $\oagg$: both states have aggregate $3$, hence lie in one fibre, and $R(\st{12})\cap R(\st{21})=\{1\}\cap\{2\}=\emptyset$; no first instrument exists. Under the codex ($\led$): the two singleton fibres issue $1$ and $2$; the successor set is $\{\st{21},\st{12}\}$ and must survive the global alternation from both phases, but $\st{21}\notin\Phi_1$.
\item $\st{12|31}$ and $\st{13|21}$ under the shared ledger (both observations): the first step issues $1$ for the state in $\{\st{12},\st{13}\}$ and $2$ for the other; the successor sets $\{\st{21},\st{22}\}$ and $\{\st{22},\st{12}\}$ must lie in $\Phi_1\cap\Phi_2$, which they do not. Under the free protocol each state keeps its own thread and both pairs are viable; under the private register they are viable as well (the threads never share a fibre with opposite phase).
\item $\st{13|31}$ under $\oagg$ (both protocols) and under the private codex: both states have aggregate $4$ and share a fibre with $R(\st{13})\cap R(\st{31})=\emptyset$ under $\oagg$; under $\ofull$ with the private register the two threads land on $\st{22}$ after one step with opposite phases, which is a coherence violation. Under the shared ledger at full observation the pair is viable: the mixed first step $(1,2)$ sends both states to $\st{22}\in\Phi_1\cap\Phi_2$.
\end{itemize}

\begin{table}[H]
\centering\footnotesize
\caption{Master table: the $36$ pairs of safe states and their membership ($\bullet$) in the kernels. Columns: institutional, aggregate, codex and full under the shared ledger; institutional and codex under the private register ($\priv$); codex under the optimistic ledger ($\opt$); decentralised unanimity. Pairs through the corner $\st{11}$ are in no kernel.}
\label{tab:master}
\begin{tabular}{lcccccccc}
\toprule
pair & inst & agg & codex & full & inst$^{\priv}$ & codex$^{\priv}$ & codex$^{\opt}$ & dec \\
\midrule
$\st{11|12}$ & $\cdot$ & $\cdot$ & $\cdot$ & $\cdot$ & $\cdot$ & $\cdot$ & $\cdot$ & $\cdot$ \\
$\st{11|13}$ & $\cdot$ & $\cdot$ & $\cdot$ & $\cdot$ & $\cdot$ & $\cdot$ & $\cdot$ & $\cdot$ \\
$\st{11|21}$ & $\cdot$ & $\cdot$ & $\cdot$ & $\cdot$ & $\cdot$ & $\cdot$ & $\cdot$ & $\cdot$ \\
$\st{11|22}$ & $\cdot$ & $\cdot$ & $\cdot$ & $\cdot$ & $\cdot$ & $\cdot$ & $\cdot$ & $\cdot$ \\
$\st{11|23}$ & $\cdot$ & $\cdot$ & $\cdot$ & $\cdot$ & $\cdot$ & $\cdot$ & $\cdot$ & $\cdot$ \\
$\st{11|31}$ & $\cdot$ & $\cdot$ & $\cdot$ & $\cdot$ & $\cdot$ & $\cdot$ & $\cdot$ & $\cdot$ \\
$\st{11|32}$ & $\cdot$ & $\cdot$ & $\cdot$ & $\cdot$ & $\cdot$ & $\cdot$ & $\cdot$ & $\cdot$ \\
$\st{11|33}$ & $\cdot$ & $\cdot$ & $\cdot$ & $\cdot$ & $\cdot$ & $\cdot$ & $\cdot$ & $\cdot$ \\
$\st{12|13}$ & $\bullet$ & $\bullet$ & $\bullet$ & $\bullet$ & $\bullet$ & $\bullet$ & $\bullet$ & $\cdot$ \\
$\st{12|21}$ & $\cdot$ & $\cdot$ & $\cdot$ & $\bullet$ & $\cdot$ & $\bullet$ & $\cdot$ & $\bullet$ \\
$\st{12|22}$ & $\bullet$ & $\bullet$ & $\bullet$ & $\bullet$ & $\bullet$ & $\bullet$ & $\bullet$ & $\cdot$ \\
$\st{12|23}$ & $\bullet$ & $\bullet$ & $\bullet$ & $\bullet$ & $\bullet$ & $\bullet$ & $\bullet$ & $\bullet$ \\
$\st{12|31}$ & $\cdot$ & $\bullet$ & $\cdot$ & $\bullet$ & $\bullet$ & $\bullet$ & $\bullet$ & $\cdot$ \\
$\st{12|32}$ & $\bullet$ & $\bullet$ & $\bullet$ & $\bullet$ & $\bullet$ & $\bullet$ & $\bullet$ & $\bullet$ \\
$\st{12|33}$ & $\bullet$ & $\bullet$ & $\bullet$ & $\bullet$ & $\bullet$ & $\bullet$ & $\bullet$ & $\cdot$ \\
$\st{13|21}$ & $\cdot$ & $\bullet$ & $\cdot$ & $\bullet$ & $\bullet$ & $\bullet$ & $\bullet$ & $\cdot$ \\
$\st{13|22}$ & $\bullet$ & $\bullet$ & $\bullet$ & $\bullet$ & $\bullet$ & $\bullet$ & $\bullet$ & $\bullet$ \\
$\st{13|23}$ & $\bullet$ & $\bullet$ & $\bullet$ & $\bullet$ & $\bullet$ & $\bullet$ & $\bullet$ & $\cdot$ \\
\bottomrule
\end{tabular}\quad
\begin{tabular}{lcccccccc}
\toprule
pair & inst & agg & codex & full & inst$^{\priv}$ & codex$^{\priv}$ & codex$^{\opt}$ & dec \\
\midrule
$\st{13|31}$ & $\cdot$ & $\cdot$ & $\bullet$ & $\bullet$ & $\cdot$ & $\cdot$ & $\bullet$ & $\bullet$ \\
$\st{13|32}$ & $\bullet$ & $\bullet$ & $\bullet$ & $\bullet$ & $\bullet$ & $\bullet$ & $\bullet$ & $\cdot$ \\
$\st{13|33}$ & $\bullet$ & $\bullet$ & $\bullet$ & $\bullet$ & $\bullet$ & $\bullet$ & $\bullet$ & $\bullet$ \\
$\st{21|22}$ & $\bullet$ & $\bullet$ & $\bullet$ & $\bullet$ & $\bullet$ & $\bullet$ & $\bullet$ & $\cdot$ \\
$\st{21|23}$ & $\bullet$ & $\bullet$ & $\bullet$ & $\bullet$ & $\bullet$ & $\bullet$ & $\bullet$ & $\bullet$ \\
$\st{21|31}$ & $\bullet$ & $\bullet$ & $\bullet$ & $\bullet$ & $\bullet$ & $\bullet$ & $\bullet$ & $\cdot$ \\
$\st{21|32}$ & $\bullet$ & $\bullet$ & $\bullet$ & $\bullet$ & $\bullet$ & $\bullet$ & $\bullet$ & $\bullet$ \\
$\st{21|33}$ & $\bullet$ & $\bullet$ & $\bullet$ & $\bullet$ & $\bullet$ & $\bullet$ & $\bullet$ & $\cdot$ \\
$\st{22|23}$ & $\bullet$ & $\bullet$ & $\bullet$ & $\bullet$ & $\bullet$ & $\bullet$ & $\bullet$ & $\cdot$ \\
$\st{22|31}$ & $\bullet$ & $\bullet$ & $\bullet$ & $\bullet$ & $\bullet$ & $\bullet$ & $\bullet$ & $\bullet$ \\
$\st{22|32}$ & $\bullet$ & $\bullet$ & $\bullet$ & $\bullet$ & $\bullet$ & $\bullet$ & $\bullet$ & $\cdot$ \\
$\st{22|33}$ & $\bullet$ & $\bullet$ & $\bullet$ & $\bullet$ & $\bullet$ & $\bullet$ & $\bullet$ & $\bullet$ \\
$\st{23|31}$ & $\bullet$ & $\bullet$ & $\bullet$ & $\bullet$ & $\bullet$ & $\bullet$ & $\bullet$ & $\cdot$ \\
$\st{23|32}$ & $\bullet$ & $\bullet$ & $\bullet$ & $\bullet$ & $\bullet$ & $\bullet$ & $\bullet$ & $\bullet$ \\
$\st{23|33}$ & $\bullet$ & $\bullet$ & $\bullet$ & $\bullet$ & $\bullet$ & $\bullet$ & $\bullet$ & $\cdot$ \\
$\st{31|32}$ & $\bullet$ & $\bullet$ & $\bullet$ & $\bullet$ & $\bullet$ & $\bullet$ & $\bullet$ & $\cdot$ \\
$\st{31|33}$ & $\bullet$ & $\bullet$ & $\bullet$ & $\bullet$ & $\bullet$ & $\bullet$ & $\bullet$ & $\bullet$ \\
$\st{32|33}$ & $\bullet$ & $\bullet$ & $\bullet$ & $\bullet$ & $\bullet$ & $\bullet$ & $\bullet$ & $\cdot$ \\
\bottomrule
\end{tabular}
\end{table}
\section{The whole-universe census and the register conventions}\label{si:universe}

Table~\ref{tab:universe} gives the viable beliefs by size and the stabilisation index of the horizon family on the whole node set; Table~\ref{tab:maximal} lists the inclusion-maximal viable beliefs. All kernels are closed under sub-beliefs (checked setwise), so each is the down-closure of its maximal members.

\begin{table}[H]
\centering\small
\caption{Viable beliefs by size for each class and convention (columns $1$--$8$ are belief sizes), with the total and the number $h^{\ast}$ of applications of $\Phi$ after which the horizon family on the whole node set is stationary ($\calK=\calW_{h^{\ast}}$).}
\label{tab:universe}
\begin{tabular}{lrrrrrrrrrr}
\toprule
class & 1 & 2 & 3 & 4 & 5 & 6 & 7 & 8 & total & $h^{\ast}$ \\
\midrule
full & 8 & 28 & 56 & 70 & 56 & 28 & 8 & 1 & 255 & 1 \\
codex ($\led$) & 8 & 25 & 39 & 32 & 13 & 2 & 0 & 0 & 119 & 2 \\
aggregate & 8 & 26 & 44 & 41 & 20 & 4 & 0 & 0 & 143 & 1 \\
institutional ($\led$) & 8 & 24 & 36 & 29 & 12 & 2 & 0 & 0 & 111 & 2 \\
codex ($\priv$) & 8 & 27 & 50 & 55 & 36 & 13 & 2 & 0 & 191 & 2 \\
institutional ($\priv$) & 8 & 26 & 44 & 41 & 20 & 4 & 0 & 0 & 143 & 1 \\
codex ($\opt$) & 8 & 27 & 50 & 55 & 36 & 13 & 2 & 0 & 191 & 2 \\
institutional ($\opt$) & 8 & 26 & 42 & 35 & 14 & 2 & 0 & 0 & 127 & 2 \\
decentralised & 8 & 12 & 8 & 2 & 0 & 0 & 0 & 0 & 30 & --- \\
\bottomrule
\end{tabular}
\end{table}
\begin{table}[H]
\centering\footnotesize
\caption{Inclusion-maximal viable beliefs (the antichains that generate the kernels by sub-belief closure).}
\label{tab:maximal}
\begin{tabular}{lp{0.78\textwidth}}
\toprule
class & maximal beliefs \\
\midrule
full & $\st{12|13|21|22|23|31|32|33}$ \\
codex ($\led$) & $\st{13|23|31|32|33}$, $\st{12|13|22|23|32|33}$, $\st{21|22|23|31|32|33}$ \\
aggregate & $\st{12|13|22|23|32|33}$, $\st{12|22|23|31|32|33}$, $\st{13|21|22|23|32|33}$, $\st{21|22|23|31|32|33}$ \\
institutional ($\led$) & $\st{12|13|22|23|32|33}$, $\st{21|22|23|31|32|33}$ \\
codex ($\priv$) & $\st{12|13|21|22|23|32|33}$, $\st{12|21|22|23|31|32|33}$ \\
institutional ($\priv$) & $\st{12|13|22|23|32|33}$, $\st{12|22|23|31|32|33}$, $\st{13|21|22|23|32|33}$, $\st{21|22|23|31|32|33}$ \\
codex ($\opt$) & $\st{12|13|22|23|31|32|33}$, $\st{13|21|22|23|31|32|33}$ \\
institutional ($\opt$) & $\st{12|23|31|32|33}$, $\st{13|21|23|32|33}$, $\st{12|13|22|23|32|33}$, $\st{21|22|23|31|32|33}$ \\
blind & $\st{12|13|22|23|32|33}$, $\st{21|22|23|31|32|33}$ \\
decentralised & $\st{12|21|23|32}$, $\st{13|22|31|33}$ \\
\bottomrule
\end{tabular}
\end{table}

\paragraph{The join of the middle classes.}
On all beliefs, $\Kfull\setminus(\Kagg\cup\Kcod)$ has $104$ members: the $64$ beliefs containing $\{\st{12},\st{21}\}$ and the $56$ beliefs containing $\{\st{13},\st{31}\}$ together with at least one of $\st{12},\st{21},\st{22}$, with $16$ beliefs in both groups; its minimal members are $\st{12|21}$, $\st{12|13|31}$, $\st{13|21|31}$ and $\st{13|22|31}$. The eight beliefs containing $\{\st{13},\st{31}\}$ that are codex-viable are $\{\st{13},\st{31}\}\cup S$ with $S\subseteq\{\st{23},\st{32},\st{33}\}$. $\Kagg\setminus\Kcod$ has $32$ members with minimal members $\st{12|31}$ and $\st{13|21}$; $\Kcod\setminus\Kagg$ has $8$ members with minimal member $\st{13|31}$.

\paragraph{Conventions not used in the main text.}
Two further conventions were reported in earlier editions: the ledger records the instrument of the \emph{first} or of the \emph{last} fibre in a fixed order of the fibres (by their smallest state). Both give $(25,26,26,28)$ on the pair population, they disagree with each other as sets, and neither coincides with the shared-ledger kernels $(24,26,25,28)$; this is why the shared ledger quantifies over all issued instruments instead of a fixed thread. The optimistic ledger gives $(26,26,27,28)$ on pairs and $(127,143,191,255)$ on beliefs; its codex kernel has the same size as the private-register codex kernel ($191$) but is a different set (its maximal beliefs are $\st{12|13|22|23|31|32|33}$ and $\st{13|21|22|23|31|32|33}$ against $\st{12|13|21|22|23|32|33}$ and $\st{12|21|22|23|31|32|33}$). Admitting the idle instrument $0$ as a third option leaves all four shared-ledger kernels unchanged on all $511$ beliefs of safe states: it is dominated.

\paragraph{Closed forms.}
The closed forms of Theorem~\mref{thm:closed} were implemented without the game engine (alternating trajectories and, for the private register, the coherence check over $24$ steps) and agree with the engine on all $511$ beliefs of safe states for the institutional and codex classes under both conventions. The blind free kernel is $(2^{\Phi_1}\cup 2^{\Phi_2})\setminus\{\emptyset\}$ with $111$ members and equals the shared-ledger institutional kernel setwise.
\section{Stationary policies}\label{si:stationary}

A deterministic stationary policy fixes the closed loop, so a belief is sustained by it exactly when each of its states is; Table~\ref{tab:stationary} reports the unions over policy classes of the sustained beliefs (Proposition~\mref{prop:stationary}). The four largest sets sustained by single alternation policies at aggregate observation are the four maximal beliefs of the aggregate kernel, $\st{12|13|22|23|32|33}$, $\st{12|22|23|31|32|33}$, $\st{13|21|22|23|32|33}$ and $\st{21|22|23|31|32|33}$; they correspond to the first instruments $(1,1,\cdot,\cdot)$, $(1,2,\cdot,\cdot)$, $(2,1,\cdot,\cdot)$ and $(2,2,\cdot,\cdot)$ on the aggregate cells $3$ and $4$ (cells $5$ and $6$ admit either). The memoryless aggregate policies sustain no state because the aggregate is non-increasing along every safe trajectory, so a memoryless aggregate policy eventually repeats one instrument, and a repeated instrument drives the other stock to $0$ within three steps. Enumerating the $2^{27}$ full-observation policies with memory is unnecessary: the $512$ memoryless ones already realise the full kernel.

\begin{table}[H]
\centering\small
\caption{Stationary-policy unions: number of deterministic policies, number of beliefs (of $255$) sustained by at least one of them, and the largest number of states sustained by a single policy.}
\label{tab:stationary}
\begin{tabular}{llrrr}
\toprule
observation & policy class & policies & beliefs & max states \\
\midrule
$\oagg$ & memoryless, $\oagg(V)\to A$ & 32 & 0 & 0 \\
$\oagg$ & one period of instrument memory & 32768 & 143 & 6 \\
$\oagg$ & alternation (first instrument per cell) & 32 & 143 & 6 \\
$\ofull$ & memoryless, $V\to A$ & 512 & 255 & 8 \\
\bottomrule
\end{tabular}
\end{table}
\section{Decentralised unanimity: law tables}\label{si:dec}

Table~\ref{tab:sustainers} gives, for each safe state, the number of law pairs (of $256$) that sustain it as a singleton; by Theorem~\mref{thm:factor} a belief is sustained if and only if these sets intersect. The $30$ sustained beliefs are the nonempty subsets of $\{\st{12},\st{21},\st{23},\st{32}\}$ and of $\{\st{13},\st{22},\st{31},\st{33}\}$ (two $4$-sets: $2\cdot 15=30$), so the sustained pairs are the $12$ pairs inside the two maximal beliefs and the family $\{\calS(x)\}$ has the Helly property on this instance. Table~\ref{tab:laws} lists the $16$ law pairs sustaining $\st{12|21}$ and Table~\ref{tab:authority} the authority census. The $256$ law pairs fall into $64$ behavioural classes (votes at $z=0$ are never consulted on safe orbits); the $16$ sustaining pairs of $\st{12|21}$ are $4$ behavioural classes, each extended by the $4$ unconsulted vote pairs. Fixing $l_1(1)=1$ leaves all $16$ sustaining pairs of $\st{12|21}$, fixing $l_2(1)=1$ leaves none.

\begin{table}[H]
\centering\small
\caption{Number of law pairs (of $256$) sustaining each safe state as a singleton.}
\label{tab:sustainers}
\begin{tabular}{lr}
\toprule
state & $|\calS(x)|$ \\
\midrule
$\st{11}$ & 0 \\
$\st{12}$ & 16 \\
$\st{13}$ & 20 \\
$\st{21}$ & 16 \\
$\st{22}$ & 32 \\
$\st{23}$ & 24 \\
$\st{31}$ & 20 \\
$\st{32}$ & 24 \\
$\st{33}$ & 16 \\
\bottomrule
\end{tabular}
\end{table}
\begin{table}[H]
\centering\small
\caption{The $16$ law pairs sustaining $\st{12|21}$: the product $L_1\times L_2$ (rows $L_1$, columns $L_2$; laws written as the votes at $z=0,1,2,3$).}
\label{tab:laws}
\begin{tabular}{lcccc}
\toprule
$l_1\backslash l_2$ & 1211 & 1212 & 2211 & 2212 \\
\midrule
1121 & (1121,1211) & (1121,1212) & (1121,2211) & (1121,2212) \\
1122 & (1122,1211) & (1122,1212) & (1122,2211) & (1122,2212) \\
2121 & (2121,1211) & (2121,1212) & (2121,2211) & (2121,2212) \\
2122 & (2122,1211) & (2122,1212) & (2122,2211) & (2122,2212) \\
\bottomrule
\end{tabular}
\end{table}
\begin{table}[H]
\centering\footnotesize
\caption{Authority census: the pair kernel when one vote at the floor is fixed, and when a whole law is fixed to a constant.}
\label{tab:authority}
\begin{tabular}{lrp{0.62\textwidth}}
\toprule
restriction & pairs & sustained pairs \\
\midrule
$l_1(1)=1$ (owner) & 12 & $\st{12|21}$, $\st{12|23}$, $\st{12|32}$, $\st{13|22}$, $\st{13|31}$, $\st{13|33}$, $\st{21|23}$, $\st{21|32}$, $\st{22|31}$, $\st{22|33}$, $\st{23|32}$, $\st{31|33}$ \\
$l_1(1)=2$ (against) & 4 & $\st{22|31}$, $\st{22|33}$, $\st{23|32}$, $\st{31|33}$ \\
$l_2(1)=1$ (against) & 4 & $\st{13|22}$, $\st{13|33}$, $\st{22|33}$, $\st{23|32}$ \\
$l_2(1)=2$ (owner) & 12 & $\st{12|21}$, $\st{12|23}$, $\st{12|32}$, $\st{13|22}$, $\st{13|31}$, $\st{13|33}$, $\st{21|23}$, $\st{21|32}$, $\st{22|31}$, $\st{22|33}$, $\st{23|32}$, $\st{31|33}$ \\
$l_1\equiv 1$ & 0 & --- \\
$l_1\equiv 2$ & 0 & --- \\
$l_2\equiv 1$ & 0 & --- \\
$l_2\equiv 2$ & 0 & --- \\
\bottomrule
\end{tabular}
\end{table}
\section{The monitoring-design census}\label{si:design}

All $B_8=4{,}140$ set partitions of the eight viable states were enumerated and the free-protocol kernel of each computed by the bitmask implementation of the belief game (greatest fixed point on the $255$ beliefs). Table~\ref{tab:design_sizes} gives the distribution of kernel sizes, Table~\ref{tab:coarsest} the sixteen coarsest adequate designs. The count of adequate designs by the closed form \meqref{eq:count} is $151+372+372+799=1{,}694$: for the independent-set partitions $\pi$ of the critical four states $\{\st{12},\st{13},\st{21},\st{31}\}$ with $|\pi|=2,3,3,4$ blocks, $\sum_{m=0}^{4}\binom4m B_m|\pi|^{4-m}$ equals $16+32+48+40+15=151$ ($|\pi|=2$), $81+108+108+60+15=372$ ($|\pi|=3$) and $256+256+192+80+15=799$ ($|\pi|=4$). Dynamic adequacy coincides with one-step adequacy on all $4{,}140$ partitions, as Theorem~\mref{thm:dynamic} requires. The minimum kernel $111$ is attained by exactly the $52$ partitions in which one cell contains all four critical states (the number of partitions of the remaining four states into blocks that may or may not join that cell: $\sum_m\binom4m B_m=1+4+12+20+15=52$), and every such kernel equals the blind kernel. The aggregate partition $\{\st{12},\st{21}\}\,|\,\{\st{13},\st{22},\st{31}\}\,|\,\{\st{23},\st{32}\}\,|\,\{\st{33}\}$ has kernel $143$ with four maximal beliefs of size $6$ and $26$ viable pairs. The horizon family stabilises after $0$, $1$ or $2$ applications of $\Phi$ for $1{,}694$, $1{,}404$ and $1{,}042$ partitions respectively. Under the shared-ledger no-repeat protocol the closed form of Theorem~\mref{thm:closed}(a) was evaluated for every partition: no partition sustains the whole viable set and the largest kernel is $119$ (the codex kernel, attained at the finest partition).

\begin{table}[H]
\centering\small
\caption{Distribution of kernel sizes over the $4{,}140$ monitoring partitions ($18$ sizes, $74$ distinct kernels).}
\label{tab:design_sizes}
\begin{tabular}{rr}
\toprule
kernel size & partitions \\
\midrule
111 & 52 \\
123 & 20 \\
129 & 10 \\
131 & 54 \\
133 & 17 \\
135 & 204 \\
137 & 12 \\
139 & 24 \\
143 & 217 \\
155 & 102 \\
159 & 348 \\
167 & 210 \\
173 & 16 \\
175 & 172 \\
179 & 39 \\
183 & 54 \\
191 & 895 \\
255 & 1694 \\
\bottomrule
\end{tabular}
\end{table}
\begin{table}[H]
\centering\footnotesize
\caption{The $16$ coarsest (two-cell) dynamically adequate monitoring designs.}
\label{tab:coarsest}
\begin{tabular}{l}
\toprule
design \\
\midrule
$\st{12|13}$ $\mid$ $\st{21|22|23|31|32|33}$ \\
$\st{12|13|22}$ $\mid$ $\st{21|23|31|32|33}$ \\
$\st{12|13|23}$ $\mid$ $\st{21|22|31|32|33}$ \\
$\st{12|13|22|23}$ $\mid$ $\st{21|31|32|33}$ \\
$\st{12|13|22|23|32|33}$ $\mid$ $\st{21|31}$ \\
$\st{12|13|23|32|33}$ $\mid$ $\st{21|22|31}$ \\
$\st{12|13|22|32|33}$ $\mid$ $\st{21|23|31}$ \\
$\st{12|13|32|33}$ $\mid$ $\st{21|22|23|31}$ \\
$\st{12|13|32}$ $\mid$ $\st{21|22|23|31|33}$ \\
$\st{12|13|22|32}$ $\mid$ $\st{21|23|31|33}$ \\
$\st{12|13|23|32}$ $\mid$ $\st{21|22|31|33}$ \\
$\st{12|13|22|23|32}$ $\mid$ $\st{21|31|33}$ \\
$\st{12|13|22|23|33}$ $\mid$ $\st{21|31|32}$ \\
$\st{12|13|23|33}$ $\mid$ $\st{21|22|31|32}$ \\
$\st{12|13|22|33}$ $\mid$ $\st{21|23|31|32}$ \\
$\st{12|13|33}$ $\mid$ $\st{21|22|23|31|32}$ \\
\bottomrule
\end{tabular}
\end{table}
\section{The four-state target}\label{si:target}

Table~\ref{tab:target} lists the $B_4=15$ partitions of $\{\st{12},\st{21},\st{22},\st{31}\}$ with their one-step adequacy. The incompatibility graph has the two edges $\st{12}-\st{21}$ and $\st{12}-\st{31}$; $5$ partitions merge the first edge, $5$ the second and $2$ both, so $15-(5+5-2)=7$ are adequate; the two adequate two-cell partitions are the two proper $2$-colourings of the path $\st{21}-\st{12}-\st{31}$ with the isolated $\st{22}$ attached to either side. The maximal common-action sets are $\{\st{12},\st{22}\}$ and $\{\st{21},\st{22},\st{31}\}$. The Helly numbers of the families of all nonempty subsets of a two- and a three-element instrument set are $2$ and $3$ (Theorem~\mref{thm:adequacy}(b)), computed by enumerating inclusion-minimal subfamilies with empty intersection.

\begin{table}[H]
\centering\small
\caption{The $15$ partitions of the four-state target $\{\st{12},\st{21},\st{22},\st{31}\}$ and their one-step adequacy ($\bullet$: every cell has a common safe instrument).}
\label{tab:target}
\begin{tabular}{lcc}
\toprule
partition & cells & adequate \\
\midrule
$\st{12}$ $\mid$ $\st{21}$ $\mid$ $\st{22}$ $\mid$ $\st{31}$ & 4 & $\bullet$ \\
$\st{12}$ $\mid$ $\st{21}$ $\mid$ $\st{22|31}$ & 3 & $\bullet$ \\
$\st{12}$ $\mid$ $\st{21|22}$ $\mid$ $\st{31}$ & 3 & $\bullet$ \\
$\st{12}$ $\mid$ $\st{21|22|31}$ & 2 & $\bullet$ \\
$\st{12}$ $\mid$ $\st{21|31}$ $\mid$ $\st{22}$ & 3 & $\bullet$ \\
$\st{12|21}$ $\mid$ $\st{22}$ $\mid$ $\st{31}$ & 3 & $\cdot$ \\
$\st{12|21}$ $\mid$ $\st{22|31}$ & 2 & $\cdot$ \\
$\st{12|21|22}$ $\mid$ $\st{31}$ & 2 & $\cdot$ \\
$\st{12|21|22|31}$ & 1 & $\cdot$ \\
$\st{12|21|31}$ $\mid$ $\st{22}$ & 2 & $\cdot$ \\
$\st{12|22}$ $\mid$ $\st{21}$ $\mid$ $\st{31}$ & 3 & $\bullet$ \\
$\st{12|22}$ $\mid$ $\st{21|31}$ & 2 & $\bullet$ \\
$\st{12|22|31}$ $\mid$ $\st{21}$ & 2 & $\cdot$ \\
$\st{12|31}$ $\mid$ $\st{21}$ $\mid$ $\st{22}$ & 3 & $\cdot$ \\
$\st{12|31}$ $\mid$ $\st{21|22}$ & 2 & $\cdot$ \\
\bottomrule
\end{tabular}
\end{table}
\section{The timing grids}\label{si:timing}

Table~\ref{tab:timing} is the $93$-cell grid of Corollary~\mref{cor:grid}: $33$ viable cells under free switching with unit damage ($D_T=T$) and $43$ under alternation with damages $(1,\tfrac12)$ ($D_1,D_2,D_3=1,\tfrac32,\tfrac52$; $21$, $16$ and $6$ viable cells for $T=1,2,3$). The strict form $z_0-1>T$ of the free identity fails exactly at $(2.0,1)$, $(3.0,2)$, $(4.0,3)$; the strict form against the integer count $\lfloor z_0-1\rfloor>T$ fails on $21$ cells, one in each row with $z_0\ge 2$. On the three-regime graph without self-loops and damages $(1,\tfrac12,\tfrac14)$ the max-plus recursion and path enumeration agree for $T\le 6$ with $D_T=1,\tfrac32,\tfrac52,3,4,\tfrac92$.

\begin{table}[H]
\centering\footnotesize
\caption{The $93$-cell timing grid: for each initial level $z_0$ (slack $z_0-1$, integer count $\lfloor z_0-1\rfloor$) and review period $T$, viability under free switching with unit damage (f) and under alternation with damages $(1,\tfrac12)$ (a).}
\label{tab:timing}
\begin{tabular}{rrcccccc}
\toprule
$z_0$ & $\lfloor z_0-1\rfloor$ & $T{=}1$ f & a & $T{=}2$ f & a & $T{=}3$ f & a \\
\midrule
1.0 & 0 & $\cdot$ & $\cdot$ & $\cdot$ & $\cdot$ & $\cdot$ & $\cdot$ \\
1.1 & 0 & $\cdot$ & $\cdot$ & $\cdot$ & $\cdot$ & $\cdot$ & $\cdot$ \\
1.2 & 0 & $\cdot$ & $\cdot$ & $\cdot$ & $\cdot$ & $\cdot$ & $\cdot$ \\
1.3 & 0 & $\cdot$ & $\cdot$ & $\cdot$ & $\cdot$ & $\cdot$ & $\cdot$ \\
1.4 & 0 & $\cdot$ & $\cdot$ & $\cdot$ & $\cdot$ & $\cdot$ & $\cdot$ \\
1.5 & 0 & $\cdot$ & $\cdot$ & $\cdot$ & $\cdot$ & $\cdot$ & $\cdot$ \\
1.6 & 0 & $\cdot$ & $\cdot$ & $\cdot$ & $\cdot$ & $\cdot$ & $\cdot$ \\
1.7 & 0 & $\cdot$ & $\cdot$ & $\cdot$ & $\cdot$ & $\cdot$ & $\cdot$ \\
1.8 & 0 & $\cdot$ & $\cdot$ & $\cdot$ & $\cdot$ & $\cdot$ & $\cdot$ \\
1.9 & 0 & $\cdot$ & $\cdot$ & $\cdot$ & $\cdot$ & $\cdot$ & $\cdot$ \\
2.0 & 1 & $\bullet$ & $\bullet$ & $\cdot$ & $\cdot$ & $\cdot$ & $\cdot$ \\
2.1 & 1 & $\bullet$ & $\bullet$ & $\cdot$ & $\cdot$ & $\cdot$ & $\cdot$ \\
2.2 & 1 & $\bullet$ & $\bullet$ & $\cdot$ & $\cdot$ & $\cdot$ & $\cdot$ \\
2.3 & 1 & $\bullet$ & $\bullet$ & $\cdot$ & $\cdot$ & $\cdot$ & $\cdot$ \\
2.4 & 1 & $\bullet$ & $\bullet$ & $\cdot$ & $\cdot$ & $\cdot$ & $\cdot$ \\
2.5 & 1 & $\bullet$ & $\bullet$ & $\cdot$ & $\bullet$ & $\cdot$ & $\cdot$ \\
2.6 & 1 & $\bullet$ & $\bullet$ & $\cdot$ & $\bullet$ & $\cdot$ & $\cdot$ \\
2.7 & 1 & $\bullet$ & $\bullet$ & $\cdot$ & $\bullet$ & $\cdot$ & $\cdot$ \\
2.8 & 1 & $\bullet$ & $\bullet$ & $\cdot$ & $\bullet$ & $\cdot$ & $\cdot$ \\
2.9 & 1 & $\bullet$ & $\bullet$ & $\cdot$ & $\bullet$ & $\cdot$ & $\cdot$ \\
3.0 & 2 & $\bullet$ & $\bullet$ & $\bullet$ & $\bullet$ & $\cdot$ & $\cdot$ \\
3.1 & 2 & $\bullet$ & $\bullet$ & $\bullet$ & $\bullet$ & $\cdot$ & $\cdot$ \\
3.2 & 2 & $\bullet$ & $\bullet$ & $\bullet$ & $\bullet$ & $\cdot$ & $\cdot$ \\
3.3 & 2 & $\bullet$ & $\bullet$ & $\bullet$ & $\bullet$ & $\cdot$ & $\cdot$ \\
3.4 & 2 & $\bullet$ & $\bullet$ & $\bullet$ & $\bullet$ & $\cdot$ & $\cdot$ \\
3.5 & 2 & $\bullet$ & $\bullet$ & $\bullet$ & $\bullet$ & $\cdot$ & $\bullet$ \\
3.6 & 2 & $\bullet$ & $\bullet$ & $\bullet$ & $\bullet$ & $\cdot$ & $\bullet$ \\
3.7 & 2 & $\bullet$ & $\bullet$ & $\bullet$ & $\bullet$ & $\cdot$ & $\bullet$ \\
3.8 & 2 & $\bullet$ & $\bullet$ & $\bullet$ & $\bullet$ & $\cdot$ & $\bullet$ \\
3.9 & 2 & $\bullet$ & $\bullet$ & $\bullet$ & $\bullet$ & $\cdot$ & $\bullet$ \\
4.0 & 3 & $\bullet$ & $\bullet$ & $\bullet$ & $\bullet$ & $\bullet$ & $\bullet$ \\
\bottomrule
\end{tabular}
\end{table}
\section{Scaling to larger caps}\label{si:family}

The worked system with cap $c$ has states $\{0,\dots,c\}^2$, the same instruments and the same floor. Table~\ref{tab:family} gives the pair census for $c=3,4,5$: in every case the full kernel is the set of corner-free pairs, the institutional kernel under the private register equals the aggregate kernel, and all inclusions of Corollary~\mref{cor:order} are strict, with the aggregate and codex kernels incomparable. The computation uses the generic game engine on the reachable node graph of the pair population.

\begin{table}[H]
\centering\small
\caption{Scaling of the pair census to larger caps: safe states, pairs of safe states, pair kernels of the four classes under the shared ledger, the institutional and codex pair kernels under the private register, and the number of corner-free pairs (which equals the full kernel in every case).}
\label{tab:family}
\begin{tabular}{rrrrrrrrrr}
\toprule
cap & safe & pairs & inst & agg & codex & full & inst$^{\priv}$ & codex$^{\priv}$ & corner-free \\
\midrule
3 & 9 & 36 & 24 & 26 & 25 & 28 & 26 & 27 & 28 \\
4 & 16 & 120 & 98 & 102 & 100 & 105 & 102 & 104 & 105 \\
5 & 25 & 300 & 266 & 272 & 269 & 276 & 272 & 275 & 276 \\
\bottomrule
\end{tabular}
\end{table}
\section{The continuous fibre}\label{si:continuous}

Table~\ref{tab:bench} tabulates the benchmark fibre of Proposition~\mref{prop:bench} and Table~\ref{tab:drift} the drift audit of Proposition~\mref{prop:drift}. The duality values of Proposition~\mref{prop:duality} are computed on the rational grids $u\in\{0,\tfrac1{100},\dots,1\}$ and $\lambda\in\{0,\tfrac1{100},\dots,1\}$: primal and dual $-\tfrac1{10}$ at $u=\tfrac12$ and $\lambda=(\tfrac12,\tfrac12)$ for the two-floor instance; primal $-1$ and dual $0$ on the nonconvex control set $\{0,1\}$ with payoff rows $(-1,+1)$ and $(+1,-1)$; primal and dual $+\tfrac1{10}$ at $u=\tfrac3{10}$ for the feasible contrast. The Helly-tight triple $\{u_1\ge 1,u_2\ge 1,u_1+u_2\le 1\}$ is pairwise feasible on the half-integer grid and jointly infeasible. The two floors $u-\tfrac12\ge 0$ and $-u-\tfrac12\ge 0$ on $[-1,1]$ are each satisfiable and jointly unsatisfiable, with $\max_u\min_j f_j(u)=-\tfrac12$.

\begin{table}[H]
\centering\small
\caption{The benchmark fibre: caps, cap sum and status against the demand $u_1+u_2\ge 2$; the margin is the dual-weighted deficiency $\mu(Y)=(Y-27/5)/10$ at $\lambda=(\tfrac12,\tfrac12)$.}
\label{tab:bench}
\begin{tabular}{lllll}
\toprule
$Y$ & $c_1(Y)$ & $c_2(Y)$ & $c_1+c_2$ & status \\
\midrule
$4$ & $\tfrac{13}{10}$ & $\tfrac{49}{50}$ & $\tfrac{57}{25}$ & feasible \\
$\tfrac{9}{2}$ & $\tfrac{5}{4}$ & $\tfrac{93}{100}$ & $\tfrac{109}{50}$ & feasible \\
$5$ & $\tfrac{6}{5}$ & $\tfrac{22}{25}$ & $\tfrac{52}{25}$ & feasible \\
$\tfrac{27}{5}$ & $\tfrac{29}{25}$ & $\tfrac{21}{25}$ & $2$ & feasible \\
$\tfrac{11}{2}$ & $\tfrac{23}{20}$ & $\tfrac{83}{100}$ & $\tfrac{99}{50}$ & obstructed, margin $\tfrac{1}{100}$ \\
$6$ & $\tfrac{11}{10}$ & $\tfrac{39}{50}$ & $\tfrac{47}{25}$ & obstructed, margin $\tfrac{3}{50}$ \\
$7$ & $1$ & $\tfrac{17}{25}$ & $\tfrac{42}{25}$ & obstructed, margin $\tfrac{4}{25}$ \\
\bottomrule
\end{tabular}
\end{table}
\begin{table}[H]
\centering\small
\caption{Certainty-equivalent drift under $g(s)=s^2$ with lead $\tfrac1{10}$: naive drift $s/5+1/100$ and corrected drift.}
\label{tab:drift}
\begin{tabular}{lll}
\toprule
$s$ & naive drift & corrected drift \\
\midrule
$1$ & $\tfrac{21}{100}$ & $0$ \\
$\tfrac{5}{4}$ & $\tfrac{13}{50}$ & $0$ \\
$\tfrac{3}{2}$ & $\tfrac{31}{100}$ & $0$ \\
$\tfrac{7}{4}$ & $\tfrac{9}{25}$ & $0$ \\
$2$ & $\tfrac{41}{100}$ & $0$ \\
\bottomrule
\end{tabular}
\end{table}
\section{Verification record}\label{si:verification}

\paragraph{Running the verifier.}
\texttt{python3 paper11c\_worked\_systems\_audit\_v5\_verification.py} (in the directory containing \texttt{work/}) recomputes every object from \texttt{work/ws\_core.py}, evaluates the \NCHECKS{} checks of Table~\ref{tab:checks}, compares the numbers quoted in the compiled manuscript and supplement with the recomputed values when the PDFs are present, and exits with status $0$ only if everything passes. The design census (all $4{,}140$ partitions, free protocol and shared ledger) accounts for most of the running time.

\begin{longtable}{lp{0.86\textwidth}}
\caption{The named checks of the verifier and their status in the shipped run.}\label{tab:checks}\\
\toprule
status & check \\
\midrule
\endfirsthead
\toprule
status & check \\
\midrule
\endhead
\bottomrule
\endfoot
PASS & S1 state grid 16 states, 9 safe, corner 11 has no safe instrument, eight other safe states have R as in eq. (R) \\
PASS & S2 aggregate z1+z2 is non-increasing under instruments 1 and 2 on all safe states (and falls only when the raised stock is at its cap) \\
PASS & S3 36 pairs of safe states, 28 corner-free; 511 nonempty beliefs of safe states, 255 of viable states \\
PASS & K1 shared-ledger pair kernels (inst, agg, codex, full) = (24, 26, 25, 28) \\
PASS & K2 shared-ledger belief kernels on the 511 beliefs = (111, 143, 119, 255), all contained in the corner-free beliefs \\
PASS & K3 full kernel = all nonempty subsets of the eight viable states; on pairs = the 28 corner-free pairs \\
PASS & K4 monotonicity inclusions inst $<$= agg $<$= full and inst $<$= codex $<$= full hold on all beliefs and all are strict \\
PASS & K5 aggregate and codex kernels incomparable; witnesses 12$|$31, 13$|$21 in agg \textbackslash\{\} codex and 13$|$31 in codex \textbackslash\{\} agg \\
PASS & K6 meet of the middle classes is the institutional kernel on pairs and on all beliefs \\
PASS & K7 join of the middle classes misses the full kernel by exactly 12$|$21 on pairs and by 104 beliefs (64 containing 12$|$21, the rest containing 13$|$31 plus one of 12, 21, 22); minimal members 12$|$21, 12$|$13$|$31, 13$|$21$|$31, 13$|$22$|$31 \\
PASS & K8 every kernel (all classes, all conventions) is closed under nonempty sub-beliefs \\
PASS & K9 maximal beliefs: inst 2x6, agg 4x6, codex 2x6+1x5, full 1x8 (shared ledger) \\
PASS & K10 horizon family on the whole node set stabilises after one application for the free classes and two for the ledger classes \\
PASS & K11 pair kernels are the size-2 members of the belief kernels (consistency of the two populations) \\
PASS & K12 admitting the idle instrument 0 leaves all four shared-ledger kernels unchanged \\
PASS & A1 alternation-safe sets Phi\_1 = \{12,13,22,23,32,33\}, Phi\_2 = \{21,22,23,31,32,33\}; union = viable states, intersection = \{22,23,32,33\} \\
PASS & A2 every alternating trajectory in Phi\_a is 2-periodic after at most one step \\
PASS & A3 shared-ledger closed form reproduces the institutional and codex kernels on all 511 beliefs \\
PASS & A4 private-register closed form (phase coherence) reproduces the private institutional and codex kernels on all 511 beliefs \\
PASS & A5 blind kernel = (2\^{}Phi\_1 U 2\^{}Phi\_2) minus empty, 111 beliefs, equal to the shared-ledger institutional kernel \\
PASS & A6 codex minus institutional (shared ledger) = the 8 beliefs \{13,31\} U S, S subset of \{23,32,33\} \\
PASS & A7 private-register kernels (inst, agg, codex, full) = (143, 143, 191, 255) on beliefs and (26, 26, 27, 28) on pairs; private institutional = aggregate setwise \\
PASS & A8 shared ledger and private register incomparable at full observation: 13$|$31 ledger-viable and private-dead, 12$|$21 the reverse \\
PASS & A9 optimistic-ledger kernels = (127, 143, 191, 255) on beliefs, (26, 26, 27, 28) on pairs; contain the shared-ledger kernels; optimistic codex differs setwise from private codex \\
PASS & A10 fixed-thread conventions (first / last fibre) give (25, 26, 26, 28) on pairs and disagree with each other setwise \\
PASS & P1 32 memoryless aggregate policies sustain no state \\
PASS & P2 32 alternation policies realise exactly the aggregate kernel (143 beliefs; four sustained sets of size 6 = the maximal beliefs) \\
PASS & P3 all 32768 one-period-memory aggregate policies realise the same set as the 32 alternation policies \\
PASS & P4 512 memoryless full-observation policies realise the full kernel; the displayed policy sustains all eight viable states \\
PASS & D1 16 laws, 256 law pairs, 64 behavioural classes \\
PASS & D2 statewise factoring: the factored kernel equals the direct orbit computation on all 511 beliefs \\
PASS & D3 decentralised kernel: 12 pairs, 30 beliefs (8+12+8+2), maximal beliefs 12$|$21$|$23$|$32 and 13$|$22$|$31$|$33; 16 of 28 coordinated-viable pairs lost; contains 12$|$21, misses 12$|$13 \\
PASS & D4 sustained pairs are exactly the pairs inside the two maximal beliefs (Helly property of the sustainer family on this instance) \\
PASS & D5 the 16 law pairs sustaining 12$|$21 form the product L1 x L2 (4 behavioural classes x 4 unconsulted votes) \\
PASS & D6 authority census: l1(1)=1 and l2(1)=2 leave 12 pairs; l2(1)=1 leaves \{13$|$22,13$|$33,22$|$33,23$|$32\}; l1(1)=2 its swap image; constant laws leave 0 \\
PASS & D7 fixing l1(1)=1 leaves 16 sustaining pairs for 12$|$21, fixing l2(1)=1 leaves 0 \\
PASS & M1 incompatibility graph on the eight viable states = K\_\{2,2\} between \{12,13\} and \{21,31\} (4 edges) \\
PASS & M2 four-state target: 15 partitions, 7 adequate = 15-(5+5-2); edges 12-21, 12-31; two adequate two-cell partitions; maximal common-action sets 12$|$22 and 21$|$22$|$31 \\
PASS & M3 Helly numbers of the two- and three-instrument universes are 2 and 3 \\
PASS & M4 design census: 4140 partitions, 1694 dynamically adequate, dynamic adequacy = one-step adequacy on every partition \\
PASS & M5 closed-form count 151 + 2*372 + 799 = 1694 agrees with the census \\
PASS & M6 16 coarsest (two-cell) adequate designs; 74 distinct kernels; 18 kernel sizes; minimum 111 attained by exactly the 52 partitions merging the four critical states \\
PASS & M7 maximal-belief sizes 6/7/8 for 927/1519/1694 partitions; pair kernels 24/26/27/28 for 52/906/1488/1694; aggregate partition: kernel 143, 4 maximal beliefs, 26 pairs, not one-step adequate \\
PASS & M8 under the shared ledger no partition is adequate and the largest kernel over all partitions is 119 \\
PASS & R1 frozen regimes: pinning identity, all 9 states viable, 36 pairs each \\
PASS & R2 mode-blind kernel = the eight corner-free states (28 pairs), equal to the full-information kernel; corner lost \\
PASS & R3 mode-blind horizon family stabilises after one application (W\_1 = W\_2 = W\_3) \\
PASS & T1 timing grid: 93 cells, 33 viable under free switching, 43 under alternation (21, 16, 6) \\
PASS & T2 strict form fails exactly at (2,1), (3,2), (4,3); floored strict form fails on 21 cells \\
PASS & T3 alternation worst damage 1, 3/2, 5/2, 3, 4, 9/2 = ceil(T/2) + floor(T/2)/2; max-plus recursion equals path enumeration on the 3-regime graph \\
PASS & C1 benchmark crossover Y* = 27/5; margin law (Y-27/5)/10; Y=6 caps 11/10, 39/50, sum 47/25, margin 3/50; Y=5 witness (6/5, 4/5) \\
PASS & C2 radius: per-cap increase below 3/50 leaves the Y=6 obstruction, 3/50 removes it \\
PASS & C3 two-floor duality: primal = dual = -1/10 at u = 1/2, lambda = (1/2, 1/2) \\
PASS & C4 nonconvex gap: primal -1 against dual 0; feasible contrast +1/10 on both sides at u = 3/10 \\
PASS & C5 Helly-tight triple pairwise feasible, jointly infeasible \\
PASS & C6 drift: naive s/5 + 1/100 (21/100 ... 41/100), corrected 0, sign law \\
PASS & C7 floors: each satisfiable, no common instrument, value -1/2 \\
PASS & F1 cap 4: 120 pairs, kernels (98, 102, 100, 105), full = corner-free, private institutional = aggregate \\
PASS & F2 cap 5: 300 pairs, kernels (266, 272, 269, 276), full = corner-free, private institutional = aggregate \\
PASS & J1 results.json agrees with the recomputation on the headline counts \\
PASS & L1 Lean modules present with theorems Wk\_mono\_obs, master\_action\_axis, master\_memory\_axis, master\_monotone and no sorry \\
PASS & X1 main text quotes the recomputed headline numbers (111, 143, 119, 255, 4,140, 1,694, 74 distinct kernels, 27/5, 3/50, -1/10, 12 of the 28, 16 coarsest, 32 alternation policies, 33 and 43, 93-cell) \\
PASS & X2 main text states the check count (64 named checks) consistently with this script \\
PASS & X3 supplement quotes the recomputed numbers (104 members, 151+372+372+799, 4 behavioural classes, 52 partitions, (25, 26, 26, 28), (127,143,191,255)) \\
\end{longtable}

\paragraph{Lean modules.}
Table~\ref{tab:lean} lists the declarations of the two Lean modules shipped with the article; they import the \texttt{Prelude} and \texttt{P1\_Obstruction} modules of the repository's \texttt{Formalizations} development and formalise the three axes of Theorem~\mref{thm:mono} as horizon-indexed inclusions. The verifier checks the presence of the declarations in the shipped sources.

\begin{table}[H]
\centering\small
\caption{Declarations of the Lean modules shipped with the article.}
\label{tab:lean}
\begin{tabular}{lll}
\toprule
module & kind & name \\
\midrule
\texttt{WS\_WorkedSystems.lean} & def & \texttt{refineObs} \\
\texttt{WS\_WorkedSystems.lean} & theorem & \texttt{Wk\_mono\_obs} \\
\texttt{WS\_WorkedSystems.lean} & theorem & \texttt{master\_action\_axis} \\
\texttt{WS\_WorkedSystems.lean} & theorem & \texttt{master\_memory\_axis} \\
\texttt{WS\_WorkedSystems\_v2.lean} & theorem & \texttt{master\_monotone} \\
\texttt{WS\_WorkedSystems\_v2.lean} & theorem & \texttt{monotone\_action\_axis} \\
\bottomrule
\end{tabular}
\end{table}
\section{Symbols}\label{si:symbols}

\begin{table}[H]
\centering\small
\begin{tabular}{@{}lp{0.70\textwidth}@{}}
\toprule
symbol & meaning \\
\midrule
$X$, $V$ & state grid $\{0,\dots,3\}^2$ and safe set $\{z_1\ge 1, z_2\ge 1\}$ (nine states) \\
$A=\{1,2\}$, $0$ & instruments; idle instrument \\
$F_u$ & transition of instrument $u$; $F_u(C)$ image of a set \\
$R(x)$, $R(C)$ & safe instruments of a state, common safe instruments of a set \\
$\st{z_1z_2}$, $\st{x|y}$ & state and belief labels \\
$o$, $\ofull$, $\oagg$, $\oblind$ & observation map; full, aggregate, blind \\
$B_y$ & fibre of the belief $B$ at observation $y$ \\
$\free$, $\nr$ & free and no-repeat protocols \\
$\led$, $\priv$, $\opt$ & shared ledger, private register, optimistic ledger \\
$\calN$, $\Phi$, $\calW_h$, $\calK$ & nodes, operator, horizon family, kernel of the belief game \\
$K(o,P)$ & viable beliefs under observation $o$ and protocol $P$ \\
$\Kfull,\Kcod,\Kagg,\Kinst,\Kdec$ & kernels of the full, codex, aggregate, institutional and decentralised classes \\
$\Vstar$, $\Rstar$, $\Gstar$ & viable singletons, instruments keeping a state in $\Vstar$, incompatibility graph on $\Vstar$ \\
$\Phi_a$ & states surviving the alternating schedule that starts with instrument $a$ \\
$I(a)$ & set of instruments issued by a profile \\
$l_i$, $T_{l_1,l_2}$, $\calS(x)$ & agency laws, unanimity closed loop, law pairs sustaining $x$ \\
$G_T$, $N(G_T)$, $B_m$ & incompatibility graph of a target, number of independent-set partitions, Bell numbers \\
$F^{r}_u$, $D_T$ & regime dynamics; worst cumulative damage over $T$ periods \\
$c_i(Y)$, $Y^{\ast}$, $\mu(Y)$ & caps, crossover and dual margin of the benchmark fibre \\
\bottomrule
\end{tabular}
\end{table}

\end{document}
